\chapter{Ablation Study and Analysis}
\label{sec:analysis}
\section{Candidate Size Analysis}

\begin{table}[t]
\centering
\small
\begin{tabular}{cccc}
\hline
\textbf{Candidate Size $k$} &
\textbf{Accuracy} &
\textbf{R@$k$} &
\textbf{Input Tokens} \\
\hline
1  & 0.47 & 0.65 & 1.30M \\
5  & 0.57 & 0.83 & 3.92M \\
10 & \textbf{0.60} & 0.91 & 6.54M \\
20 & \textbf{0.60} & 0.95 & 12.35M \\
\hline
\end{tabular}
\caption{Effect of candidate size $k$ on end-to-end accuracy, table retrieval recall, and LLM input-token usage.}
\label{tab:k_ablation}
\end{table}

Increasing the candidate size improves gold-table coverage and provides the reasoner with more opportunities to access the correct table.
Table~\ref{tab:k_ablation} reports end-to-end accuracy, table retrieval recall, and LLM input-token usage at different candidate sizes.
Increasing $k$ from $1$ to $5$ improves retrieval recall from $0.65$ to $0.83$ and end-to-end accuracy from $0.47$ to $0.57$.
At $k=10$, retrieval recall further increases to $0.91$, while accuracy reaches $0.60$.

The benefit of increasing candidate size exhibits diminishing returns.
Increasing $k$ allows more candidate tables, including the gold table, to remain accessible to the reasoner, but also introduces additional non-gold tables and increases the amount of information processed during reasoning.
Notably, increasing $k$ from $10$ to $20$ does not further improve accuracy, which remains at $0.60$, while input-token usage increases substantially from $6.54$M to $12.35$M.
The results therefore reveal a trade-off among gold-table coverage, reasoning difficulty, and inference cost when selecting the candidate size.

\section{Table Size Analysis}

\begin{table*}[t]
\centering
\small
\begin{tabular}{lcccccccc}
\hline
& \multicolumn{2}{c}{\textbf{Closed-domain}}
& \multicolumn{2}{c}{\textbf{Ours}}
& \multicolumn{2}{c}{\textbf{OpenTab+r}}
& \multicolumn{2}{c}{\textbf{CRAFT}} \\
\textbf{LLM}
& $<1$M & $\geq1$M
& $<1$M & $\geq1$M
& $<1$M & $\geq1$M
& $<1$M & $\geq1$M \\
\hline
GPT-5.6-Sol
& 0.88 & 0.80
& \textbf{0.62} & \textbf{0.55}
& 0.39 & 0.29
& 0.27 & 0.26 \\

Gemini-3.7-Flash
& 0.81 & 0.77
& \textbf{0.56} & \textbf{0.54}
& 0.39 & 0.31
& 0.29 & 0.26 \\

Qwen3.5-9B
& 0.78 & 0.72
& \textbf{0.63} & \textbf{0.57}
& 0.23 & 0.25
& 0.20 & 0.18 \\

Qwen3.5-4B
& 0.78 & 0.65
& \textbf{0.60} & \textbf{0.55}
& 0.22 & 0.23
& 0.22 & 0.21 \\
\hline
\end{tabular}
\caption{Accuracy on DataBench stratified by the serialized token size of the gold table. $<1$M and $\geq1$M denote gold tables containing fewer than one million and at least one million tokens, respectively. Closed-domain evaluation directly provides the gold table to the reasoner.}
\label{tab:table_size_analysis}
\end{table*}

We further analyze end-to-end performance with respect to the size of the gold table on DataBench.
We divide the evaluation queries into two subsets according to whether the fully serialized gold table contains fewer than one million tokens or at least one million tokens.
Table~\ref{tab:table_size_analysis} reports the accuracy of each method on the two subsets across all evaluated LLM backbones.

Overall performance generally decreases as the gold-table size increases.
Our approach achieves lower accuracy on the $\geq1$M-token subset across all four LLM backbones.
Despite the decrease, our approach remains the best-performing end-to-end OLTQA method on the $\geq1$M-token subset across all evaluated LLMs.

The closed-domain results provide clearer evidence of the difficulty introduced by large tables.
Even when the gold table is directly provided, accuracy consistently decreases on the $\geq1$M-token subset across all four LLM backbones.

In contrast, OpenTab+r and CRAFT do not exhibit a consistent size-dependent degradation across all LLMs.
Their performance remains low even on the $<1$M-token subset, indicating that their weak end-to-end performance on DataBench is not confined to tables containing at least one million tokens.
As shown in Table~\ref{tab:dataset_comparison}, DataBench contains substantially larger tables on average than the benchmarks previously used to evaluate these methods.
Therefore, tables in the $<1$M-token subset may still be considerably larger than those encountered in their original evaluation settings, despite falling below our operational large-table threshold.


\section{Token Efficiency Analysis}
\begin{table}[t]
\centering
\small
\begin{tabular}{lc}
\hline
\textbf{Method} & \textbf{Input Tokens} \\
\hline
OpenTab+r & 12.46M \\
CRAFT     & 6.75M \\
\hline
Ours      & \textbf{6.54M} \\
\hline
\end{tabular}
\caption{Total generative LLM input-token usage of the evaluated OLTQA pipelines on DataBench.}
\label{tab:token_efficiency}
\end{table}

Our method remains competitive in LLM input-token usage despite performing interactive multi-table reasoning.
Table~\ref{tab:token_efficiency} compares the total generative LLM input tokens consumed by the evaluated OLTQA pipelines.
OpenTab~\cite{opentab} and OpenTab+r consume $19.31$M and $12.46$M input tokens, respectively, while CRAFT~\cite{craft} consumes $6.75$M.
Our method consumes $6.54$M input tokens, remaining competitive with the evaluated baselines despite allowing the reasoner to interact with multiple candidate tables.

This difference may partly reflect the additional LLM-based processing required by the baseline pipelines.
OpenTab uses an LLM-based coder for SQL generation, while CRAFT uses LLM-based table enrichment before the reasoning stage.
In contrast, our pipeline does not require separate SQL-generation or table-enrichment stages.


\section{Interactive Table Access Analysis}

\begin{table}[t]
\centering
\small
\begin{tabular}{lc}
\hline
\textbf{Method} & \textbf{Accuracy} \\
\hline
Ours (Interactive) & \textbf{0.61} \\
Ours (Static)      & 0.25 \\
\hline
\end{tabular}
\caption{Ablation of interactive table access on DataBench using Qwen3.5-9B. The static variant disables access to additional table information during reasoning while retaining the same initial table context.}
\label{tab:interactive_ablation}
\end{table}

We analyze the importance of interactive table access for reasoning over large tables by comparing our full method with a static variant.
The static variant uses the same initial table context as our full method but disables programmatic access to additional table information during reasoning.
This comparison isolates the contribution of interactive table access while keeping the initial table context unchanged.

As shown in Table~\ref{tab:interactive_ablation}, disabling interactive table access reduces accuracy substantially from $0.61$ to $0.25$.
The large performance decrease indicates that the initial table context alone is insufficient for many queries and that access to additional table information during reasoning is important for successful OLTQA over large tables.
Interactive table access allows the reasoner to obtain additional information when the initially constructed context does not contain sufficient information for answering the query.


\section{Retriever Ablation Study}
\label{sec:retriever_ablation}

To evaluate the contribution of the hybrid retrieval strategy, we compare BM25 sparse retrieval, dense retrieval using \texttt{sentence-transformers/all-mpnet-base-v2}, and their hybrid combination while keeping the schema-level table representation fixed.

\begin{table}[t]
\centering
\caption{Ablation study of retrieval algorithms on the DataBench test set using schema-level table representations.}
\label{tab:retriever_ablation}
\begin{tabular}{lcccc}
\toprule
\textbf{Retriever} &
\textbf{R@1} &
\textbf{R@3} &
\textbf{R@5} &
\textbf{R@10} \\
\midrule
BM25
& 0.56 & 0.66 & 0.72 & 0.80 \\
Dense (all-mpnet-base-v2)
& 0.56 & 0.72 & \textbf{0.83} & 0.90 \\
Hybrid (BM25 + all-mpnet-base-v2)
& \textbf{0.65} & \textbf{0.77} & \textbf{0.83} & \textbf{0.91} \\
\bottomrule
\end{tabular}
\end{table}

Table~\ref{tab:retriever_ablation} compares the retrieval performance of different retrieval algorithms under the same schema-level table representation.
Hybrid retrieval achieves the highest Recall@1 of 0.65, compared with 0.56 for both BM25 and dense retrieval. It also achieves the highest Recall@3 and Recall@10, while matching dense retrieval at Recall@5.

These results indicate that combining sparse and dense retrieval improves top-ranked table retrieval performance, although the gains become smaller at larger candidate depths.

\paragraph{Dense Embedding Model Comparison.}

We additionally evaluate a more recent dense embedding model, \texttt{jinaai/jina-embeddings-v4}~\footnote{\url{https://huggingface.co/jinaai/jina-embeddings-v4}}, to examine the sensitivity of our hybrid retriever to the choice of dense encoder.
We compare it with the original \texttt{sentence-transformers/all-mpnet-base-v2} under the same hybrid retrieval framework.
This comparison is conducted as a supplementary analysis on the DataBench test set, rather than as part of development-set model selection.

\begin{table}[t]
\centering
\caption{Comparison of dense embedding models in our hybrid retriever on the DataBench test set. Both configurations use BM25 and schema-level table representations.}
\label{tab:embedding_comparison}
\begin{tabular}{lcccc}
\toprule
\textbf{Dense Encoder} &
\textbf{R@1} &
\textbf{R@3} &
\textbf{R@5} &
\textbf{R@10} \\
\midrule
all-mpnet-base-v2
& \textbf{0.65} & \textbf{0.77} & \textbf{0.83} & \textbf{0.91} \\
jina-embeddings-v4
& 0.53 & 0.66 & 0.72 & 0.81 \\
\bottomrule
\end{tabular}
\end{table}

As shown in Table~\ref{tab:embedding_comparison}, the hybrid retriever using \texttt{sentence-transformers/all-mpnet-base-v2} achieves higher recall across all evaluated candidate depths.
In particular, Recall@1 decreases from 0.65 to 0.53 when replacing \texttt{sentence-transformers/all-mpnet-base-v2} with \texttt{jinaai/jina-embeddings-v4}.

These results indicate that adopting a more recent embedding model does not necessarily improve retrieval performance in our evaluation setting.
Since this comparison was conducted as a supplementary analysis after the main experiments, the original embedding configuration remains unchanged, and no additional end-to-end experiments are conducted with \texttt{jinaai/jina-embeddings-v4}.


\section{Retrieval--Reasoning Combination Analysis}
\label{sec:retrieval_reasoning_combination}

Although CRAFT achieves table retrieval recall comparable to our method, its end-to-end QA accuracy is substantially lower. To further investigate this performance gap, we evaluate an additional configuration that combines CRAFT's retriever with our interactive multi-table reasoner.

Specifically, we compare three configurations: the original CRAFT pipeline, our proposed pipeline, and CRAFT's retriever combined with our reasoner (CRAFT+Ours). In CRAFT+Ours, the top-$10$ candidate tables retrieved by CRAFT are passed to our reasoner, which constructs its own initial table context and performs interactive reasoning. All configurations are evaluated on the DataBench test set using Qwen3.5-9B, with a candidate size of $k=10$. All other experimental settings remain consistent with the main experiments.


\begin{table}[t]
\centering
\caption{Retrieval and end-to-end QA performance of different retrieval--reasoning combinations on the DataBench test set using Qwen3.5-9B ($k=10$).}
\label{tab:retrieval_reasoning_combination}
\small
\begin{tabular}{lllccccc}
\toprule
\textbf{Method} &
\textbf{Retriever} &
\textbf{Reasoner} &
\multicolumn{4}{c}{\textbf{Retrieval Recall}} &
\textbf{E2E Acc.} \\
\cmidrule(lr){4-7}
& & &
\textbf{R@1} &
\textbf{R@3} &
\textbf{R@5} &
\textbf{R@10} & \\
\midrule
CRAFT
& CRAFT & CRAFT
& 0.56 & 0.78 & 0.86 & 0.92 & 0.20 \\
Ours
& Ours & Ours
& 0.65 & 0.77 & 0.83 & 0.91 & 0.61 \\
CRAFT+Ours
& CRAFT & Ours
& 0.56 & 0.78 & 0.86 & 0.92 & -- \\
\bottomrule
\end{tabular}
\end{table}


Table~\ref{tab:retrieval_reasoning_combination} compares the retrieval recall and end-to-end QA accuracy of the three retrieval--reasoning configurations. While CRAFT and our method achieve similar Recall@10, their end-to-end performance differs substantially. The CRAFT+Ours configuration allows us to examine whether this performance gap persists when using the same interactive multi-table reasoner with a different retrieval component.

% TODO: Add the CRAFT+Ours accuracy and discuss the observed result.
